\cotefile{58}
\scanpage{11}
que $Y'\cap V(f_{\deleted{\unsure{1}}2})\cap\cdots\cap V(f_{\unsure{k|h}})\inserted{=X'\cap V(f_1)\cdots\cap V(f_k)}$ est disconnexe,
donc on est \unsure{ramené} au

\uline{Corollaire 1} Sous les conditions

\noindent\begin{tabular}{@{}|p{.89\linewidth}@{}}
$\unsure{a_k)|h')} ,\, b_k)$ soient $f_1\cdots f_k\in rA$, alors
$Y'=V(f_1)\cap\cdots V(f_k)\cap X'$ est connexe.
\end{tabular}

Or pour prouver ceci, on
est ramené au cas où $A$ est
\uline{complet}, ce que nous supposons
\unsure{dorénavant}.

Pour $k=1$, c'est le lemme 2.

Supposons $k\geq2$, et le th.
démontré pour les $k'<k$. \inserted{\unsure{Par réduct.,|Par récurr.,}} On est
ramené à prouver le th pour
$k$, et $l=1$. Procédons comme
tantôt, s'il était faux pour $Y'=V(f_1)\cap X'$,
on \unsure{disconnecterait|déconnecterait|déconnectant} \unsure{par} $\deleted{Z'}$ par $Z'$ de
dim $<k-1$, donc pour $\exists\,f_2\cdots f_k\in rA$
tels que $V(f_{\unsure{2|1}})\cap\deleted{\unsure{Z|Z'}}\cap V(f_k)=\varnothing$, et que
% archive sheet number: 10
\scanpage{12}
\mspage{- 3 -}\hfill n\handwritten{°354}

\uline{L'anneau} $A$, \uline{muni de cette graduation}, \uline{est alors appelé anneau gradué de type}
$\handwritten{\Delta}$ .

PROPOSITION 1.- \uline{Si tout élément de} $\handwritten{\Delta}$ \uline{est régulier}, \uline{et si} $(A_{\handwritten{\lambda}})$ \uline{est une graduation de type}
$\handwritten{\Delta}$ \uline{compatible avec la structure d'un anneau} $A$, $A_{\unsure{o|0}}$ \uline{est un}
\uline{sous-anneau de} $A$ (\uline{et en particulier on a} $1\handwritten{\in} A_o$)

Comme $A_o A_o\handwritten{\subset} A_o$ par définition, il suffit de prouver la seconde assertion.
Soit $1=\handwritten{\sum}\limits_{\handwritten{\lambda\in\Delta}}e_{\handwritten{\lambda}}$ la décomposition de 1 en ses composantes homogènes. Si
$x\handwritten{\in}A_\mu$, on a $x=x.1=\handwritten{\sum}\limits_{\handwritten{\lambda\in\Delta}}xe_{\handwritten{\lambda}}$ ; comparant les composantes de degré $\mu$,
on obtient (puisque la relation $\mu+\handwritten{\lambda}=\mu$ entraîne $\handwritten{\lambda}=0$) $x=xe_o$. Cette
relation étant vraie pour tout élément homogène de $A$, est vraie pour tout
$x\handwritten{\in}A$ ; en particulier $1=1.e_o=e_o\handwritten{\in} A_o$ .

DEFINITION 3.- \uline{Soient} $A$ \uline{un anneau gradué de type} $\handwritten{\Delta}$ , $M$ \uline{un} $A$-\uline{module à gauche}
(resp. \uline{à droite}) ; \uline{on dit qu'une graduation} $(M_{\handwritten{\lambda}})$ \uline{de type} $\handwritten{\Delta}$ \uline{est compatible}
\uline{avec la structure de} $A$-\uline{module de} $M$ \uline{si l'on a}
\begingroup
\makeatletter\tagsleft@true\makeatother
\[
 A_{\handwritten{\lambda}} M_\mu \handwritten{\subset} M_{\handwritten{\lambda}+\mu}
 \qquad (\text{resp. }M_\mu A_{\handwritten{\lambda}}\handwritten{\subset} M_{\handwritten{\lambda}+\mu})
 \tag*{(2)}
\]
\endgroup
\uline{quels que soient} $\handwritten{\lambda}$, $\mu$ \uline{dans} $\handwritten{\Delta}$ . \uline{Le module} $M$, \uline{muni de cette graduation}, \uline{est}
\uline{alors appelé module gradué à gauche} (resp. \uline{à droite}) \uline{de type} $\handwritten{\Delta}$ .

Lorsque les éléments de $\handwritten{\Delta}$ sont réguliers, il résulte de (2) et de la
prop.1 que les $M_{\handwritten{\lambda}}$ sont des $A_o$-\uline{modules}.

Il est clair que si $A$ est un anneau gradué de type $\handwritten{\Delta}$ , le $A$-module à gauche $A_s$ (resp. le $A$-module à droite $A_d$) est gradué de type $\handwritten{\Delta}$ .

\uline{Exemples}.- 4) Sur un anneau quelconque $A$, la graduation triviale de type
est compatible avec la structure d'anneau. Si $A$ est gradué par la graduation
triviale, pour qu'une graduation $(M_{\handwritten{\lambda}})$ de type $\handwritten{\Delta}$ sur un $A$-module $M$ soit compatible avec la structure de $A$-module, il faut et il suffit que les $M_{\handwritten{\lambda}}$
soient des \uline{sous-modules} de $M$.
% archive sheet number: 11
\scanpage{13}
un tel système, $X'\cap V(f_1)\cap\cdots V(f_{\unsure{k|h}})$
est \uline{\unsure{disconnexe|disconnecté|disconnexes}}. \uline{\uuline{Or}} \deleted{\unsure{pour}} en vertu
$\deleted{\unsure{X'}}$ du lemme 3, $\exists\,\unsure{\Phi|\varphi}\subset X'$ fini
tel que $f\in rA$, $V(f)\cap\Phi=\varnothing\ \Longrightarrow$
$X'\cap V(f)$ est connexe en dim $\leq k-1$.

\deleted{On peut supposer que $\Phi$ contient les}
On peut choisir les $f_i$ tels que
$f_2$ soit tel que $V(f_2)\cap\Phi=\varnothing$, \deleted{\unsure{il}}
[il suffit que $f_2$ ne s'annule pas
sur $\Phi$ ni les \unsure{pts} \deleted{\unsure{gén.}}\inserted{maximaux} de $Z'$].

Alors $X'\cap V(f_2)$ satisfait les
hyp. $a_{k-1})$, $b_{k-1})$. En vertu
de l'hyp. de récurrence, on
doit avoir $\bigl(X'\cap V(f_2)\bigr)\cap V(f_1)\cap V(f_3)\cap\cdots\cap V(f_k)$
connexe, ce qui est \uline{absurde}.

% slanted
\margin{\uuline{Remarque} $\searrow$}
\uline{Corollaire} \uline{I}nversement, \inserted{$(k\geq1)$} supposons que

\margin{\uline{Excluons} le cas $X$ irréductible \deleted{\unsure{Excluant|Exclure} le cas $k=1$, $X$ irréductible} \unsure{et} de dim 1,}
\noindent\begin{tabular}{@{}|p{.89\linewidth}@{}}
pour \unsure{tte} suite $f_1\cdots f_k$ d'\unsure{élém\textsuperscript{ts}|élés} de
$rA$, on ait $\pi_0(X'\cap V(f_1)\cap\cdots\cap V(f_k))$\deleted{\unsure{connexe}}$\unsure{\leq 1}$.
\deleted{connexe par \unsure{dim} $\unsure{\geq k-1|{\geq k+1}}$,} Alors $\deleted{\unsure{X'}}$ \unsure{les}
comp. de $X'$ sont de dim $\geq k+1$, et $X'$ est connexe
en dim $\geq k$.
\end{tabular}
% archive sheet number: 12
\scanpage{14}
\mspage{- 2 -}\hfill n\handwritten{°354}

duation de type $\handwritten{\Delta}'$ sur $G$.

Par exemple, soient $\handwritten{\Delta}_{\unsure{o|0}}$ un monoïde commutatif ayant un élément neutre 0,
$I$ un ensemble d'indices arbitraire, $\handwritten{\Delta}=\handwritten{\Delta}_o^{(I)}$, sous-monoïde du produit
$\handwritten{\Delta}_0^I$ formé des familles $(\handwritten{\lambda}_{\handwritten{\unsure{\iota|i}}})_{\handwritten{\iota}\handwritten{\in}I}$ à support fini. Soit $(n_{\handwritten{\iota}})_{\handwritten{\iota}\handwritten{\in}I}$ une famille
quelconque d'entiers $\handwritten{\geq}0$ ; l'application $(\handwritten{\lambda}_{\handwritten{\iota}})\handwritten{\to}\handwritten{\sum}\limits_{\handwritten{\iota}\handwritten{\in}I}n_{\handwritten{\iota}}\handwritten{\lambda}_{\handwritten{\iota}}$ est un homomorphisme $\handwritten{\rho}$ de $\handwritten{\Delta}$ dans $\handwritten{\Delta}_o$ tel que $\handwritten{\rho}(0)=0$, et associe donc à toute graduation de type $\handwritten{\Delta}$ une graduation de type $\handwritten{\Delta}_o$.

3) Plus particulièrement, si dans l'exemple 2), on prend $\handwritten{\Delta}=\handwritten{\Delta}_o\handwritten{\times}\handwritten{\Delta}_o$, on
dit qu'une graduation de $\handwritten{\Delta}$ est une \uline{bigraduation} de type $\handwritten{\Delta}_o$ ; les trois
homomorphismes $(\handwritten{\lambda},\mu)\handwritten{\to}\handwritten{\lambda}$, $(\handwritten{\lambda},\mu)\handwritten{\to}\mu$, $(\handwritten{\lambda},\mu)\handwritten{\to}\handwritten{\lambda}+\mu$ associent alors
à toute bigraduation de type $\handwritten{\Delta}_o$ trois graduations de type $\handwritten{\Delta}_o$ dites respectivement \uline{première} et \uline{seconde graduation partielle}, et \uline{graduation totale}. De
façon plus explicite, si $(G_{\handwritten{\lambda}\mu})_{(\handwritten{\lambda},\mu)\handwritten{\in}\handwritten{\Delta}}$ est la bigraduation donnée, la
première graduation partielle est formée des $G'_{\handwritten{\lambda}}=\handwritten{\bigoplus}\limits_{\mu\handwritten{\in}\handwritten{\Delta}_o}G_{\handwritten{\lambda}\mu}$, la seconde
des $G''_{\handwritten{\lambda}}=\handwritten{\bigoplus}\limits_{\mu\handwritten{\in}\handwritten{\Delta}_o}G_{\mu \handwritten{\lambda}}$ et la graduation totale des $G'''_{\handwritten{\lambda}}=\handwritten{\bigoplus}\limits_{\mu+\unsure{y|\nu}=\handwritten{\lambda}}G_{\mu y}$. On
notera que l'on a $G'_{\handwritten{\lambda}}\handwritten{\cap}G''_\mu=G_{\handwritten{\lambda}\mu}$ ; inversement, si $(G'_{\handwritten{\lambda}})$ et $(G''_{\handwritten{\lambda}})$ sont deux
graduations de type $\handwritten{\Delta}_o$ sur un groupe commutatif $G$, telles que $G$ soit somme
directe des sous-groupes $G'_{\handwritten{\lambda}}\handwritten{\cap}G''_\mu=G_{\handwritten{\lambda}\mu}$, ces sous-groupes définissent une
bigraduation de type $\handwritten{\Delta}$ sur $G$, dont $(G'_{\handwritten{\lambda}})$ et $(G''_{\handwritten{\lambda}})$ sont les deux graduations
partielles.

Nous laissons au lecteur le soin de généraliser ces définitions au cas o\unsure{\handwritten{\`}{u}|\handwritten{\=}{u}}
$\handwritten{\Delta}=\handwritten{\Delta}_o^n$ .

DEFINITION 2.- \uline{Etant donnés un anneau} $A$ \uline{et une graduation} $(A_{\handwritten{\lambda}})$ \uline{de type} $\handwritten{\Delta}$
\uline{sur le groupe additif} $A$, \uline{on dit que cette graduation est compatible avec la}
\uline{structure d'anneau de} $A$ \uline{si l'on a}
\begingroup
\makeatletter\tagsleft@true\makeatother
\[
 A_{\handwritten{\lambda}}A_\mu\handwritten{\subset}A_{\handwritten{\lambda}+\mu}\qquad
 \text{\uline{quels que soient} }\handwritten{\lambda}\ \text{\uline{et} }\mu\ \text{\uline{dans} }\handwritten{\Delta}\ .
 \tag*{(1)}
\]
\endgroup
% archive sheet number: 13
\scanpage{15}
En effet, \unsure{évidemt|évident|évidemment} $X'\ne\varnothing$. Soit
$X'_0$ une composante de dim \unsure{minima|minimum},
soit $l$, \inserted{donc $l\geq0$} prouvons $l\unsure{\geq|>} k+1$. \deleted{Sinon il
existera $f_1\cdots f_l$} \deleted{Si on avait $l\leq k$,
alors $X'$ étant connexe [prendre des
$f_i=0$] on \unsure{conclurait|concluerait} que $X'=X'_0$, donc
$X'$ de dim 0, mais alors $\exists\,f_1\in rA$
tel que $Y'=X'\cap V(f)=\varnothing$, \deleted{\unsure{à fortiori}}\inserted{\unsure{que}}
$X'\cap\underbrace{V(f)\cap\cdots\cap V(f)}_{k\ \text{fois}}=\varnothing$, absurde.}
\deleted{Si on avait $l\leq k$,} il existe\deleted{\unsure{ra}}
$l$ \unsure{élémts|éléments} de $rA$ \deleted{tels que} $f_1\cdots f_l$
tels que, \inserted{\unsure{fixant} \deleted{\unsure{pour|tous}} \unsure{les|tte|toute} composantes irréductibles $X'_i$, de dim $n_i$,}
$V(f_1)\cap\cdots\cap \unsure{V|N}(f_l)\cap X'_{\deleted{\unsure{o}}i}\unsure{(=}Y'_i)$ \deleted{\unsure{soit de dim}}
\inserted{les \unsure{comp. irréd.} de $Y'_i$} \deleted{ait \unsure{tt. comp.}} de dim $n_i-l$, et \deleted{\unsure{alors ces composantes}} irréductibles
\deleted{\unsure{fermés} \unsure{et} \unsure{que} \unsure{les} \unsure{pts} \unsure{soient} \unsure{clairs}} sont \unsure{de} distinctes ($i$ variable).

\deleted{\unsure{de ttes} composantes irréd. de $X'\cap V(f_1)\cap\cdots\cap V(f_l)$.
Comme ce dernier est connexe, il
serait de dim 0 et \unsure{réduit à}
\unsure{un} pt \unsure{Si on avait} $l\leq k$,}

En particulier, pour $i=0$, $Y'_i$ est fini,
donc les pts de $Y'_i$ sont des composantes
irréductibles de dim 0. \inserted{\uuline{Si} $l\leq k$,} l'hyp.
implique que $Y'$ est connexe, donc
$X'$ serait irréductible de dim $l\leq k$
\deleted{Si $l<k$, \unsure{absurde}, \unsure{car} une \unsure{équation|section} de}
% archive sheet number: 14
\scanpage{16}
\uline{ALGEBRE}, \uline{CHAPITRE II} (3e édition)\hfill n\handwritten{°354}

\handwritten{\S} 8. \uline{Modules et anneaux gradués}.

\uline{Dans ce paragraphe}, \uline{on désignera par} $\handwritten{\Delta}$ \uline{un monoïde} (Chap.I, \handwritten{\S} 1,n\handwritten{°3}) \uline{commutatif}, \uline{noté additivement}, \uline{ayant un élément neutre noté} 0.

1. \uline{Groupes commutatifs gradués} ; \uline{anneaux et modules gradués}.

DEFINITION 1.- \uline{Etant donné un groupe commutatif} $G$, \uline{noté additivement}, \uline{on appelle graduation de type} $\handwritten{\Delta}$ \uline{sur} $G$ \uline{une famille} $(G_{\handwritten{\lambda}})_{\handwritten{\lambda}\handwritten{\in}\handwritten{\Delta}}$ \uline{de sous-groupes de} $G$, \uline{dont} $\unsure{G|G'}$ \uline{soit somme directe}. \uline{L'ensemble} $G$, \uline{muni de la structure définie par sa loi de groupe et sa graduation}, \uline{est appelé groupe} (\uline{commutatif}) \uline{gradué de type} $\handwritten{\Delta}$ .

On dit que $\handwritten{\Delta}$ est le \uline{monoïde des degrés} de $G$. On dit qu'un élément $x\handwritten{\in}G$
est \uline{homogène} s'il appartient à un des $G_{\handwritten{\lambda}}$, \uline{homogène de degré} $\handwritten{\lambda}$ si $x\handwritten{\in}G_{\handwritten{\lambda}}$ .
L'élément 0 est donc homogène de tous les degrés ; mais si $x\handwritten{\ne}0$ est homogène
il appartient à un seul des $G_{\handwritten{\lambda}}$ ; l'indice $\handwritten{\lambda}$ tel que $x\handwritten{\in}G$ est alors appelé
\uline{le degré de} $x$ (ou quelquefois \deleted{le} \uline{poids} de $x$) et se note souvent $\deg(x)$.
Tout $y\handwritten{\in}G$ s'écrit d'une seule manière comme somme $\handwritten{\sum}\limits_{\handwritten{\lambda}}y_{\handwritten{\lambda}}$ d'éléments homogènes avec $y_{\handwritten{\lambda}}\handwritten{\in}G_{\handwritten{\lambda}}$ ; on dit que $y_{\handwritten{\lambda}}$ est la \uline{composante homogène de degré}
(ou simplement la \uline{composante de}\quad \uline{degré} $\handwritten{\lambda}$) de $y$.

\uline{Exemples}.- 1) Etant donnés un monoïde commutatif $\handwritten{\Delta}$ quelconque (ayant un
élément neutre) et un groupe commutatif quelconque $G$, on définit une graduation $(G_{\handwritten{\lambda}})_{\handwritten{\lambda}\handwritten{\in}\handwritten{\Delta}}$ sur $G$ en prenant $G_{\unsure{o|0}}=G$ et $G_{\handwritten{\lambda}}=\handwritten{\{}0\handwritten{\}}$ pour $\handwritten{\lambda}\handwritten{\ne}0$ ; cette graduation est dite \uline{triviale}.

2) Soit $\handwritten{\Delta}'$ un second monoïde commutatif ayant un élément neutre, \inserted{\handwritten{et}} soit $\handwritten{\rho}$ un
homomorphisme de $\handwritten{\Delta}$ dans $\handwritten{\Delta}'$ tel que $\handwritten{\rho}(0)=0$. Soit $(G_{\handwritten{\lambda}})_{\handwritten{\lambda}\handwritten{\in}\handwritten{\Delta}}$ une graduation de type $\handwritten{\Delta}$ sur un groupe commutatif $G$ ; pour $\mu\handwritten{\in}\handwritten{\Delta}'$, soit $G'_\mu$ la
somme des $G_{\handwritten{\lambda}}$ tels que $\handwritten{\rho}(\handwritten{\lambda})=\mu$ ; il est clair que $(G'_\mu)_{\mu\handwritten{\in}\handwritten{\Delta}'}$ est une gra-
% archive sheet number: 15
\scanpage{17}
\deleted{plus \unsure{donnant|on aurait|précisément} le vide. Si $l=k$,
on aurait $\dim X'=k$ i.e. $\dim X=k+1\geq2$,
et tte famille $f_1\cdots f_k$, $Y'=X'\cap V(f_1)\cdots\cap V(f_k)$
connexe. Or \deleted{\unsure{on trouverait} \unsure{toujours}
\unsure{facilement} \unsure{un cas} où $\dim X=2$}
\unsure{car} si $l=1$}
\unsure{i.e.} $X$ irréd de dim $l+1$\deleted{\unsure{?}}.
Alors nous \unsure{excluons}
le cas $l=0$, \inserted{i.e.} i.e. \deleted{\unsure{$X$|de}} irréd. de dim 1.

Nous avons donc \unsure{montré|à montrer|montrer} que si $X$ est
irréductible de dim \inserted{$\unsure{l|k}+1$}$\geq2$, on peut trouver
\deleted{\unsure{les|l}} \unsure{élémts|des} $f_1\cdots f_l$ de $rA$ tels que
$X'\cap V(f_1)\cap\cdots\cap V(f_l)$ disconnexe. Or
il existe $f_1\cdots f_{l-1}$ tels que
$Y'=X'\cap V(f_1)\cap\cdots\cap V(f_{l-1})$ soit de dim 1,
donc \unsure{corresponde|correspond} à un $Y$ de dim 2.

Mais alors il est facile de
trouver $f\in rB$ tel que \unsure{$V(f)$|$V(f_l)$} \unsure{de|sur} \unsure{$Y$,|$Y'$}
[qui est de dim 0] ne soit pas
\unsure{réduit|vide} : \unsure{mais|car} \unsure{alors} --- \unsure{en|car}
$Y'$ \unsure{$\subset$|où} on trouve deux \unsure{éléments|distincts|courbes} !!
% archive sheet number: 16
\scanpage{18}
\mspage{- 10 -}\hfill n\handwritten{°358}

sur $K'$. Comme $W'$ engendre l'espace $K$-vectoriel $W$, l'ensemble $H$ des $h\handwritten{\in}R'$
nulles sur $W'$ est aussi l'ensemble des $h\handwritten{\in}R'$ nulles sur $W$ ; comme toute
forme $K'$-linéaire sur $V'$ s'étend de manière unique en un élément de $R'$, il
résulte du th.8 du n\handwritten{°5} du \handwritten{\S\ 4} que $W'$ est l'ensemble des $v\handwritten{\in}V'$ tels que
$h(v)=0$ pour tout $h\handwritten{\in}H$. L'intersection des noyaux des éléments de $H$ est
un sous-espace $K$-vectoriel $W_1$ de $V$ rationnel sur $K'$, et ce qui précède montre que l'on a $W_1\handwritten{\cap}V'=W'=W\handwritten{\cap}V'$, d'o\unsure{\handwritten{\`}{u}|\handwritten{\=}{u}} $W=W_1$ d'après la prop.2. C.Q.F.D.

\uline{Remarques}.- a) On démontre facilement que si $v$ est un vecteur d'un espace
$K$-vectoriel $V$ muni d'une $K'$-structure $V'$, le corps $K'(v)$ est engendré par $K'$
et les scalaires $h(v)$ pour $h$ parcourant l'ensemble des formes linéaires
rationnelles sur $K'$. Soit $W$ un sous-espace $K$-vectoriel de $V$, que nous munirons de $K'$-structure induite par celle de $V$. Comme toute forme linéaire rationnelle sur $K'$ définie sur $W$ se prolonge en une forme linéaire rationnelle
définie sur $V$, on voit que pour tout $v\handwritten{\in}W$, le corps $K'(v)$ calculé ``dans $W$ ''
est égal au corps $K'(v)$ calculé ``dans $V$ ''. Autrement dit, pour tout sous-corps $L$ de $K$ contenant $K'$, on a $W_L=V_L\handwritten{\cap}W$.

b) Lorsque $V$ est de dimension finie, son dual $V^{\handwritten{*}}$ est engendré comme espace
$K$-vectoriel à gauche par l'ensemble $R'$ des formes linéaires rationnelles sur
$K'$. D'après le lemme 2, on voit que $R'$ est alors une $K'$-structure sur $V^{\handwritten{*}}$. On
montre facilement que pour tout sous-corps $L$ de $K$ contenant $K'$, $V_L^{\handwritten{*}}=L.R'$ est
l'ensemble des $f\handwritten{\in}V^{\handwritten{*}}$ tels que $f(V_L)\handwritten{\subset}L$ ; il n'y a donc aucune \unsure{ambiguïté|ambigüité} à
parler dans ce cas d'une ``forme linéaire sur $V$ rationnelle sur $L$ ''. Il n'en
est plus de même lorsque $V$ est de dimension infinie sur $K$ (cf. exerc.\quad ).

6. \uline{Critères de rationalité}.

Si $L$ est un sous-corps de $K$ l'ensemble $K$ considéré comme espace vectoriel
\uline{à gauche} sur $L$ sera noté $ {}_L K$ ; l'ensemble des endomorphismes de cet espace
% archive sheet number: 17
\scanpage{19}
\uline{Corollaire} $X'$ normal \inserted{[\unsure{géom\textsuperscript{t}|géom.} \unsure{unibranche} suffit] connexe} de dim $\geq n$,

\noindent\begin{tabular}{@{}|p{.89\linewidth}@{}}
$Y'=X'\cap H_1\cap\cdots\cap H_{\unsure{k|h}}$, où $k\leq n-1$. Alors
\deleted{\unsure{$\pi_1$|$\pi_0$}} $Y'$ est connexe, et $\pi_1(Y')\mathrel{\text{---}}\pi_{\unsure{2|1}}(X')$
est surjectif
\end{tabular}

En effet, la connexité : \unsure{déjà|étant} \unsure{pr.|fait|acquise}.
Pour la surjectivité des $\pi_1$, il faut
prouver que si $X'_1$ est un \unsure{rev.|r.v.} étale
\uline{connexe} de $X'$, alors $Y'_1$ est connexe.
Or $X'_1$ \unsure{sera|est} \uline{irréductible}, \deleted{donc} de dim $\geq n$,
d'où le résultat \unsure{\ldots}.

\deleted{\uline{\unsure{Application|Question}}} \uline{\unsure{Conséquence|Conséquences}} Pour vérifier
l'\unsure{engendrement|engendremt|ingrédient} fini des $\pi_1(X')$
dans le cas d'un anneau local $A$,
on est ramené [par \unsure{dévissage|devt}] à le
vérifier pour les anneaux \uline{normaux}
de dim 2. Conjecture : cela \unsure{marche}
\deleted{\unsure{si le corps}} pour les anneaux locaux
\uline{complets}, à corps résiduel \unsure{alg\textsuperscript{t}|alg\textsuperscript{mt}} clos.
% archive sheet number: 18
